%
% Exemplo LaTeX de monografia UNISINOS
%
% Elaborado com base nas orientações dadas no documento
% ``GUIA PARA ELABORAÇÃO DE TRABALHOS ACADÊMICOS''
% disponível no site da biblioteca da Unisinos.
% http://www.unisinos.br/biblioteca
%
% Os elementos textuais abaixo são apresentados na ordem em que devem
% aparecer no documento.  Repare que nem todos são obrigatórios - isso
% é devidamente indicado em cada caso.
%
% Comentários abaixo colocados entre aspas (`` '') foram
% extraídos diretamente do documento da biblioteca.
%
% Este documento é de domínio público.
%

%=======================================================================
% Declarações iniciais identificando a classe de documento e
% selecionando alguns pacotes adicionais.
%
% As opções disponíveis (separe-as com vírgulas, sem espaço) são:
% - twoside: Formata o documento para impressão frente-e-verso
%   (o default é somente-frente)
% - english,brazilian,french,german,etc.: idiomas usados no documento.
%   Deve ser colocado por último o idioma principal.
%=======================================================================
\documentclass[twoside,english,brazilian]{UNISINOSmonografia}
\usepackage[utf8]{inputenc} % charset do texto (utf8, latin1, etc.)
\usepackage[T1]{fontenc} % encoding da fonte (afeta a sep. de sílabas)
\usepackage{graphicx} % comandos para gráficos e inclusão de figuras
\usepackage{bibentry} % para inserir refs. bib. no meio do texto
\usepackage{hyperref}

%=======================================================================
% Escolha do sistema para geração de referências bibliográficas.
%
% O default é usar o estilo unisinos.bst.  Comente a definição abaixo
% e descomente a linha seguinte para usar o estilo do ABNTeX (é
% necessário ter esse pacote instalado).
%
% A vantagem do unisinos.bst é que ele permite o uso de um arquivo .bib
% seguindo as orientações tradicionais do BibTeX (veja essas orientações
% em http://ctan.tug.org/tex-archive/biblio/bibtex/contrib/doc/btxdoc.pdf).
% Entretanto, o estilo não suporta algumas citações mais exóticas como
% apud.  Para isso, use o ABNTeX, mas esteja ciente de que muitas de
% suas referências serão incompatíveis com os estilos tradicionais do
% BibTeX como plain, alpha, ieeetr, entre outros.
%=======================================================================
\unisinosbst
%\usepackage[alf]{abntcite}

%=======================================================================
% Dados gerais sobre o trabalho.
%=======================================================================
\autor{Simão}{Samuel}
\titulo{Exemplo de Monografia no Padrão da UNISINOS}
\subtitulo{Versão \LaTeX}
\orientador[Orientador(a): Prof(a). Esp., Ms. (ou Me.) ou Dr.]{Knuth}{Donald E.}
%\coorientador[Prof.~Dr.]{Lamport}{Leslie}

\unidade{Unidade Acadêmica de Graduação ou Unidade Acadêmica de Educação Online}
\curso{Bacharelado em Sistemas de Informação}
\nivel{Graduação}
\natureza{Trabalho de Conclusão de Curso apresentado como requisito parcial para a obtenção do título de Bacharel em Sistemas de Informação, pelo Curso de Sistemas de Informação da Universidade do Vale do Rio dos Sinos --- UNISINOS}
\local{São Leopoldo}
\ano{2011}

% cada palavra-chave deve ser fornecida duas vezes, uma em português e
% outra no idioma estrangeiro (na verdade, em tantos idiomas quantos se
% desejar).
\palavrachave{brazilian}{UNISINOS}
\palavrachave{brazilian}{ABNT}
\palavrachave{brazilian}{\LaTeX}
\palavrachave{brazilian}{Formatação de documentos}
\palavrachave{english}{UNISINOS}
\palavrachave{english}{ABNT}
\palavrachave{english}{\LaTeX}
\palavrachave{english}{Document typesetting}

%=======================================================================
% Início do documento.
%=======================================================================
\begin{document}
\capa
\folhaderosto

%=======================================================================
% Dedicatória (opcional).
%
% O texto é normalmente colocado na parte de baixo da página, alinhado
% à direita.  Mas a formatação é basicamente livre.  Só não se escreve
% a palavra 'dedicatória'.
%=======================================================================
\begin{dedicatoria}
Aos nossos pais.\\[4ex] % quebra a linha dando um espaçamento maior
\begin{itshape} % faz o texto ficar em itálico
If I have seen farther than others,\\
it is because I stood on the shoulders of giants.\\
\end{itshape}
--- \textsc{Sir Isaac Newton} % \textsc é o "small caps"
\end{dedicatoria}

%=======================================================================
% Agradecimentos (opcional).
%=======================================================================
\begin{agradecimentos}
Obrigado!
\end{agradecimentos}

%=======================================================================
% Epígrafe (opcional).
%
% ``[...] o autor apresenta uma citação, seguida de indicação de autoria,
% relacionada com a matéria tratada no corpo do trabalho. Podem, também,
% constar epígrafes nas folhas de aberturas das seções primárias.''
%=======================================================================
\begin{epigrafe}
``\textit{Ninguém abre um livro sem que aprenda alguma coisa}''.\\
(Anônimo)
\end{epigrafe}

%=======================================================================
% Resumo em Português.
%
% A recomendação é para 150 a 500 palavras.
%=======================================================================
\begin{abstract}
Este documento apresenta orientações para uso da classe \LaTeX\ de formatação de monografias para a UNISINOS\@.  Ao mesmo tempo, ele serve como exemplo de uso da classe, demonstrando os principais comandos a serem utilizados, e outras orientações mais gerais de uso do \LaTeX.  Adicionalmente, procuramos incluir no documento algumas orientações sobre a escrita da monografia em si, reunindo dicas e recomendações que contribuem para aumentar a qualidade técnica dos trabalhos acadêmicos.  O Resumo deve conter de 150 a 500~palavras e nele não deve haver citações. Sugere-se a utilização de parágrafo único.
\end{abstract}

%=======================================================================
% Resumo em língua estrangeira (obrigatório somente para teses e
% dissertações).
%
% O idioma usado aqui deve necessariamente aparecer nos parâmetros do
% \documentclass, no início do documento.
%=======================================================================
\begin{otherlanguage}{english}
\begin{abstract}
This document presents guidelines on the use of UNISINOS's \LaTeX\ class for academic reports and dissertations.  At the same time, it serves as an example on using the class, employing the main commands and providing further general orientations on the use of \LaTeX.  In addition, we have added guidelines for the process of writing itself, collecting tips and recommendations that contribute to the technical quality enhancement of academic monographs.  The Abstract should be composed of 150 to 500~words and must not contain any citations.  It is suggested that a single paragraph be used.
\end{abstract}
\end{otherlanguage}

%=======================================================================
% Lista de Figuras (opcional).
%=======================================================================
\listoffigures

%=======================================================================
% Lista de Tabelas (opcional).
%=======================================================================
\listoftables

%=======================================================================
% Lista de Abreviaturas (opcional).
%
% Deve ser passada como parâmetro a maior das abreviaturas utilizadas.
%=======================================================================
\begin{listadeabreviaturas}{seg., segs.}
\item[ampl.] ampliado, -a
\item[atual.] atualizado, -a
\item[coord.] coordenador
\item[N.~T.] Novo Testamento
\item[seg., segs.] seguinte, -s
\end{listadeabreviaturas}

%=======================================================================
% Lista de Siglas (opcional).
%
% Deve ser passada como parâmetro a maior das siglas utilizadas.
%=======================================================================
\begin{listadesiglas}{FAPERGS}
\item[ABNT] Associação Brasileira de Normas Técnicas
\item[CAPES] Coordenação de Aperfeiçoamento de Pessoal de Nível Superior
\item[FAPERGS] Fundação de Amparo à Pesquisa do Estado do Rio Grande do Sul
\end{listadesiglas}

%=======================================================================
% Lista de Símbolos (opcional).
%
% Deve ser passado o maior (mais largo) dos símbolos utilizados.
%=======================================================================
\begin{listadesimbolos}{Ca}
\item[\textsuperscript{o}C] Graus Celsius
\item[Al] Alumínio
\item[Ca] Cálcio
\end{listadesimbolos}

%=======================================================================
% Sumário
%=======================================================================
\tableofcontents

%=======================================================================
% Introdução
%=======================================================================
\chapter{Introdução}

% as epígrafes nos capítulos são opcionais
\epigrafecap{The reasonable man adapts himself to the world; the unreasonable one persists in trying to adapt the world to himself. Therefore all progress depends on the unreasonable man.}{George Bernard Shaw}

Conforme \citetexto{Hexsel11}, a introdução tem o objetivo de ``\emph{introduzir} o material que vai ser apresentado em mais detalhe nas seções subseqüentes''. Na introdução você deve contextualizar o problema e mostrar por que vale a pena resolvê-lo. Você deve apresentar a solução proposta e mostrar o seu diferencial em relação aos trabalhos relacionados. Observe, porém, que na introdução você deve apenas tratar do O QUÊ e PORQUÊ, sem tratar do como \cite{Hexsel11}, que deve ser explicado na seção que descreve o trabalho desenvolvido.

Geralmente, a introdução tem uma estrutura similar ao resumo e deve apresentar:
\begin{itemize}
	\item \textbf{Contexto e motivação:} Aqui você deve apresentar o contexto do trabalho (área de que ele se trata) e uma motivação para trabalhar nesse assunto.
	\item \textbf{Problema:} Aqui você vai apresentar um problema, uma lacuna, observada na área e que você pretende tratar. Você deve se perguntar aqui: ``Que respostas estou disposto a responder?''. O problema deve ser definido claramente e delimitado em termos de espaço de tempo. Veja que essa parte visa alertar o leitor de que o que você está propondo é uma solução para um problema observado na área. 
	\item \textbf{Objetivos:} Aqui você deve apresentar os objetivos do seu trabalho. Tome cuidado para não confundir objetivos com atividades.   Faça a si mesmo a pergunta: ``O que pretendo alcançar com a pesquisa?''. Você pode discernir entre objetivos gerais e objetivos específicos:
	\begin{itemize}
		\item Objetivo geral --- qual o propósito da pesquisa?
		\item Objetivos específicos --- abertura do objetivo geral em outros menores (possíveis capítulos).
	\end{itemize}
	Veja abaixo um exemplo de objetivo retirado da monografia de~\citetexto{Teixeira09}:

	Com a possibilidade de acesso a base de dados XML gerada a partir do Sistema de Currículos Lattes e a necessidade de melhor reutilizar as informações existentes neste sistema, o presente trabalho tem como objetivo geral permitir o acesso do pesquisador a seus dados através de uma interface mais amigável: o padrão LaTeX. Para isto destacam-se os seguintes objetivos específicos:
	\begin{alineas}
		\item identificar e analisar o formato de especificação de currículos da Plataforma Lattes;
		\item disponibilizar uma ferramenta para a geração de uma representação de dados intermediária a partir do formato especificado;
		\item implementar a tradução dos dados colhidos em código LaTeX através da utilização da ferramenta criada;
		\item analisar os resultados obtidos e as alternativas presentes no uso da ferramenta.
	\end{alineas}
\end{itemize}





%=======================================================================
% Referencial Teórico
%=======================================================================
\chapter{Referencial teórico}

\section{Notas fiscais de consumidor eletrônicas no Brasil}

A Nota Fiscal Eletrônica (NF-e) é um documento fiscal digital, gerado e armazenado eletronicamente, utilizado para registrar operações comerciais. No Brasil, sua implantação foi iniciada em 2005 com o objetivo de modernizar os documentos fiscais e substituir gradualmente os registros em papel por documentos eletrônicos com validade jurídica. Segundo a Secretaria da fazenda do Rio Grande do Sul [\url{https://www.sefaz.rs.gov.br/ASP/SEF_ROOT/inf/SEF-NFE_Duvidas1.htm}], a NF-e a adoção da NF-e traz diversos benefícios como a redução de custos com impressão e armazenagem de documentos, a simplificação de obrigações acessórias e a maior agilidade no acesso às informações fiscais.

Uma característica importante da NF-e é a presença da chave de acesso, composta por 44 dígitos, utilizada para identificar individualmente o documento e permitir sua consulta em ambientes oficiais da administração tributária, como o site nfe.fazenda.gov.br [\url{https://www.nfe.fazenda.gov.br/portal/principal.aspx}]. Por meio dessa chave, destinatários e demais interessados autorizados podem verificar a existência e a validade da nota fiscal eletrônica.

Além da NF-e tradicional, existem diferentes modelos de documentos fiscais eletrônicos destinados a finalidades específicas. Entre eles, destacam-se a NF-e modelo 55 e a Nota Fiscal de Consumidor Eletrônica (NFC-e), modelo 65. Enquanto a NF-e é utilizada em diversas operações de circulação de mercadorias, a NFC-e é voltada às vendas realizadas diretamente ao consumidor final. Por estar relacionada ao consumo no varejo, a NFC-e é o documento de interesse deste estudo. [\url{https://www.contabilizei.com.br/contabilidade-online/tipos-de-notas-fiscais/}]


%!!!!
% Adicionar uma imagems de uma DANFE aqui???
%!!!


\section{Características da NFC-e}

A Nota Fiscal de Consumidor Eletrônica (NFC-e) documenta compras comuns do varejo, como aquelas realizadas em supermercados, farmácias e outros estabelecimentos comerciais.

Para facilitar o acesso do consumidor às informações da compra, a NFC-e possui um documento auxiliar denominado DANFE NFC-e. Esse documento consiste em uma representação simplificada da transação de venda no varejo e permite ao consumidor consultar a NFC-e eletrônica no ambiente da Secretaria da Fazenda.

De acordo com as especificações técnicas do DANFE NFC-e [\url{https://www.nfe.fazenda.gov.br/portal/exibirArquivo.aspx?conteudo=p4adPc9Q0bA\%3D}], a área destinada ao detalhamento dos produtos ou serviços adquiridos deve apresentar, no mínimo, os seguintes campos: código, descrição, quantidade, unidade de medida, valor unitário e valor total. Esses campos permitem identificar os itens adquiridos, suas quantidades e os valores correspondentes, constituindo uma fonte estruturada de dados sobre as compras realizadas pelo consumidor.

\section{Processamento de Linguagem Natural}

O Processamento de Linguagem Natural (PLN) é uma área da ciência da computação voltada ao tratamento computacional da linguagem humana. Como os computadores não operam diretamente sobre o significado das palavras da mesma forma que os seres humanos, os textos precisam ser processados e representados de modo que seus padrões possam ser analisados por algoritmos. Segundo a IBM [\url{https://www.ibm.com/br-pt/think/topics/natural-language-processing}], o PLN é um subcampo da ciência da computação que utiliza técnicas de aprendizado de máquina para permitir que computadores interpretem e se comuniquem por meio da linguagem humana.

No contexto deste trabalho, o PLN é relevante porque as descrições dos itens presentes em notas fiscais são formadas por textos em linguagem natural, frequentemente curtos, abreviados e pouco padronizados. Para que essas descrições possam ser utilizadas na classificação automática de produtos, é necessário submetê-las a etapas de tratamento e transformação que permitam aos modelos identificar padrões textuais associados a diferentes categorias

\section{Representação de textos para modelos de aprendizado de máquina}

Para que textos possam ser utilizados como entrada em modelos de aprendizado de máquina, é necessário transformá-los em uma representação estruturada e processável computacionalmente. Isso ocorre porque algoritmos de machine learning não operam diretamente sobre o texto em sua forma bruta, mas sobre atributos numéricos extraídos desses textos. Segundo Manning, Raghavan e Schütze (2008), o processamento textual envolve etapas como a definição da unidade documental, a tokenização, o pré-processamento linguístico dos tokens e a construção de um vocabulário de termos. A tokenização consiste em dividir a sequência de caracteres do texto em unidades menores, chamadas tokens, que posteriormente podem ser normalizadas e utilizadas na representação do documento.

%!!!!
% Adicionar uma imagem de texto representado para um modelo???
%!!!

Entre as formas comuns de representação textual, destacam-se o modelo bag-of-words, os n-grams e a ponderação TF-IDF. No modelo bag-of-words, o texto é representado com base nos termos que aparecem no documento, geralmente desconsiderando a ordem original das palavras. Os n-grams ampliam essa representação ao considerar sequências de uma ou mais palavras, permitindo capturar combinações frequentes de termos. Já o TF-IDF atribui pesos aos termos considerando sua frequência em um documento e sua relevância em relação ao conjunto de documentos, reduzindo a importância de palavras muito frequentes e pouco discriminativas. Dessa forma, essas técnicas permitem converter descrições textuais, como itens de notas fiscais, em atributos que podem ser utilizados por modelos de classificação ou agrupamento.


\section{Classificação de textos com Machine Learning}

A classificação de textos é uma tarefa de aprendizado de máquina que consiste em atribuir automaticamente uma classe ou categoria previamente definida a um texto, com base em seus padrões de conteúdo. Segundo a Elastic [\url{https://www.elastic.co/pt/what-is/text-classification}], esse processo analisa o texto e associa a ele o rótulo considerado mais adequado, sendo utilizado em aplicações como detecção de spam, análise de sentimentos e organização de documentos por tema. No contexto deste trabalho, essa técnica pode ser aplicada às descrições de produtos presentes em notas fiscais, de modo que cada item seja associado a uma categoria de consumo previamente estabelecida.

\section{Tipos de Aprendizado}

Métodos de aprendizado de máquina podem ser classificados, de acordo com a forma de treinamento, em supervisionados, não supervisionados e semi-supervisionados. Segundo Chapelle, Schölkopf e Zien (2006), no aprendizado supervisionado, o conjunto de treinamento é composto por pares de entrada e saída (xi, yi), de modo que o objetivo do algoritmo é aprender um mapeamento de x para y. No aprendizado não supervisionado, por outro lado, não há rótulos ou valores-alvo associados aos dados de treinamento; nesse caso, o objetivo é identificar estruturas, padrões ou relações presentes nos dados. Já o aprendizado semi-supervisionado situa-se entre essas duas abordagens, pois utiliza uma combinação de dados rotulados e não rotulados. Nesse cenário, apenas uma parte dos exemplos possui saída conhecida, enquanto os demais exemplos permanecem sem rótulo.


\section{Algoritmos de clusterização}

Segundo Manning, Raghavan e Schütze (2008), algoritmos de clusterização têm como objetivo agrupar dados em conjuntos, chamados clusters, com base em padrões de similaridade. Em uma boa clusterização, espera-se que os elementos pertencentes ao mesmo cluster sejam semelhantes entre si, enquanto elementos de clusters diferentes apresentem maior dissimilaridade. 


\begin{figure}
    \caption{Exemplo de clusterização}
    \centering
    \includegraphics[width=0.5\linewidth]{images/clustering.png}
    \label{fig:placeholder}
    \fonte{Manning, Raghavan e Schütze (2008)}
\end{figure}


Esse tipo de algoritmo é associado ao aprendizado não supervisionado, pois os dados não possuem classes previamente atribuídas por um especialista. Diferentemente da classificação, em que o modelo aprende a partir de exemplos rotulados, na clusterização a formação dos grupos é determinada pela própria distribuição dos dados. Assim, a escolha da medida de distância ou similaridade utilizada pelo algoritmo influencia diretamente os agrupamentos encontrados.


\section{Avaliação de modelos de classificação}

A avaliação de modelos de classificação busca mensurar a qualidade das previsões realizadas pelo algoritmo. Para isso, é comum analisar a relação entre os rótulos reais e os rótulos previstos, normalmente representada por uma matriz de confusão, da qual podem ser derivadas métricas como acurácia, precisão, recall e F1-score. Segundo [\url{https://www.mdpi.com/1999-4893/16/5/236}], juntamente com tempo de execução, esses são as 5 métricas mais utilizadas para avaliação de modelos.

A acurácia corresponde à proporção total de previsões corretas realizadas pelo modelo. A precisão indica, entre os itens classificados em determinada categoria, quantos realmente pertencem a ela. O recall representa, entre todos os itens que de fato pertencem a uma categoria, quantos foram corretamente identificados pelo modelo. Já o F1-score combina precisão e recall em uma única medida, sendo especialmente útil quando se deseja equilibrar os dois aspectos da avaliação.

Em problemas multiclasse, como a categorização de produtos em diferentes grupos de consumo, essas métricas podem ser calculadas individualmente para cada classe e também de forma agregada. Isso é relevante porque um modelo pode apresentar bom desempenho geral, mas ainda assim ter baixa qualidade em categorias menos representadas no conjunto de dados.











%=======================================================================
% Escrevendo o Texto
%=======================================================================
\chapter{Escrevendo o Texto}

\section{Comandos do \LaTeX}
Como regra geral, use os comandos tradicionais do \LaTeX\ para formatar seu texto.  Neste documento procuramos demonstrar os comandos mais comumente utilizados em monografias acadêmicas.

Neste capítulo apresentamos alguns exemplos de como colocar figuras e tabelas no seu texto.

\section{Ilustrações}

\subsection{Legendas}
As legendas das figuras devem se encontrar no topo da figura e não abaixo, como usualmente colocado. Abaixo da figura, é obrigatório colocar a fonte (mesmo que a figura tenha sido do próprio autor).

As legendas devem conter o tipo da ilustração (Figura, Tabela, etc), seguido de numeração simples (sem número do capítulo).

Toda figura deve ser citada no texto, como nos exemplos que seguem.

\subsection{Figuras}
A Figura~\ref{fig:escrita} ilustra as fases psicológicas da escrita da dissertação. Você vai se reconhecer no personagem. ;-)

\begin{figure}
	\caption{Fases psicológicas da escrita da dissertação}
	\label{fig:escrita}
	\centering%
	\begin{minipage}{.8\textwidth}
		\includegraphics[width=\textwidth]{escrita}
		\fonte{http://www.phdcomics.com/comics/archive.php?comicid=149}
	\end{minipage}
\end{figure}

\subsection{Tabelas}
A Tabela~\ref{tab:estacoes} é um exemplo de tabela elaborada pelo(a) próprio(a) autor(a).

\begin{table}
	\caption{Período das estações do ano no Brasil}
	\label{tab:estacoes}
	\centering%
	\begin{minipage}{.6\textwidth}
		\begin{tabular*}{\textwidth}{ll}
			\hline
			\textbf{Meses} & \textbf{Estações do Ano}\\
			\hline
			21 de março a 21 de junho & Outono\\
			21 de junho a 23 de setembro & Inverno\\
			23 de setembro a 21 de dezembro & Primavera\\
			21 de dezembro a 21 de março & Verão\\
			\hline
		\end{tabular*}
		\fonte{Elaborada pela autora.}
	\end{minipage}
\end{table}

\section{Resumo}
O resumo deve conter de 100 a 500 palavras. No resumo não deve haver citações e indica-se que essa seja a última seção do texto a ser escrita. Veja abaixo uma sugestão de organização e exemplo de resumo de \citetexto{Moro11}.

Sugestão (uma a três linhas para cada item):
\begin{itemize}
	\item Contexto geral e específico;
	\item Questão/problema sendo investigado (propósito do trabalho);
	\item Estado-da-arte (por que precisa de uma solução nova/melhor);
	\item Solução (nome da proposta, metodologia básica sem detalhes, quais características respondem as questões iniciais, interpretação dos resultados, conclusões).
\end{itemize}

Exemplo (SANTOS et al., 2008 apud \citealp{Moro11}):
\begin{quote}
CONTEXTO: A Web é abundante em páginas que armazenam  dados de forma implícita. PROBLEMA: Em muitos casos, estes dados estão presentes em textos semiestruturados sem a presença de delimitadores explícitos e organizados em uma estrutura também implícita. SOLUÇÃO: Este artigo apresenta uma nova abordagem para extração em textos semi-estruturados baseada em Modelos de Markov Ocultos (Hidden Markov Models - HMM). ESTADO-DA-ARTE e MÉTODO PROPOSTO: Ao contrário de outros trabalhos baseados em HMM, a abordagem proposta dá ênfase à extração de metadados, além dos dados propriamente ditos. Esta abordagem consiste no uso de uma estrutura aninhada de HMMs, onde um HMM principal identifica os atributos no texto e HMMs internos, um para cada atributo, identificam os dados e metadados. Os HMMs são gerados a partir de um treinamento com uma fração de amostras da base a ser extraída. RESULTADOS: Os experimentos realizados com anúncios de classificados retirados da Web mostram que o processo de extração alcança qualidade acima de 0,97 com a medida F, mesmo se esta fração de treinamento é pequena. 
\end{quote}

%=======================================================================
% Exemplos de Citações e Referências Bibliográficas
%=======================================================================
\chapter{Exemplos de Citações e Referências Bibliográficas}
\nobibliography* % para usar o \bibentry
Neste capítulo são apresentados exemplos de citações e referências bibliográficas.  Aqui é utilizado o pacote \texttt{bibentry}, que permite a inserção de referências no meio do texto (atenção para a diferença entre citações e referências).

Você vai ver que, neste exemplo, não está sendo usado o estilo de referências bibliográficas do projeto ABNTeX\footnote{http://http://sourceforge.net/projects/abntex}.  Você é completamente livre para usá-lo (veja no início do arquivo .tex como fazer isso).  Os motivos para não usar o ABNTeX neste exemplo são basicamente dois:
\begin{itemize}
	\item Para usar o ABNTeX, é necessário instalá-lo em seu sistema \TeX\ primeiro; embora não seja uma tarefa tão complicada, enxergamos como uma dificuldade a mais para o usuário iniciante.  Nosso objetivo aqui é facilitar ao aluno da UNISINOS o uso deste modelo, de modo que basta copiar os arquivos \texttt{UNISINOSmonografia.cls} e \texttt{unisinos.bst} para a pasta onde estão seus arquivos .tex;
	\item As normas da ABNT são tão complexas que, para atender a todas as variações possíveis de citações e referências, o projeto ABNTeX criou uma série de campos adicionais nas entradas do arquivo .bib.  Embora funcione para o caso ABNT, o efeito colateral de fazer isso é que o seu arquivo .bib será muitas vezes incompatível com os demais estilos tradicionais do BibTeX, como \texttt{plain}, \texttt{alpha}, \texttt{ieeetr}, entre outros.  Por exemplo, em referências a artigos publicados em conferências, o campo \texttt{organization} é usado pelo ABNTeX para definir o nome do evento.  Isso não é padrão e não será reconhecido pelos estilos tradicionais\footnote{Veja como criar seus arquivos .bib no manual do BibTeX, que pode ser encontrado em http://ctan.tug.org/tex-archive/biblio/bibtex/contrib/doc/btxdoc.pdf.}.  Considerando que um dos maiores benefícios do BibTeX é criar um arquivo .bib que pode ser reutilizado pelo resto da vida, nossa estratégia com o \texttt{unisinos.bst} foi tentar aproximar ao máximo a formatação exigida pela ABNT sem implicar na criação de arquivos .bib incompatíveis.  Isso funciona bem na grande maioria dos casos, mas não em todos.  Nesse caso, a saída é usar o ABNTeX ou então alterar manualmente o arquivo .bbl que é gerado ao rodar o comando \texttt{bibtex}.
\end{itemize}

Em caso de dúvida, siga as orientações do manual da Biblioteca \cite{Biblioteca11} e, se necessário, da norma NBR~6023 \cite{NBR6023:2002}.

\section{Citações}
As citações podem ocorrer de duas formas: com os nomes dos autores inseridos no texto ou não.  Isso implica em uma construção diferente para as frases.  Por exemplo:
\begin{itemize}
	\item Com o nome do autor inserido no texto: ``De acordo com \citetexto{Tanenbaum03}, o modelo de referência OSI foi proposto de forma tardia.''
	\item Sem inserir o autor no texto: ``O modelo de referência OSI foi proposto de forma tardia \cite{Tanenbaum03}.''
\end{itemize}

\section{Livros}
Seguem alguns exemplos de referências de livros:
\begin{itemize}
	\item \bibentry{Buford09}.
	\item Livro com indicação de edição:\\
	\bibentry{Kurose10ptbr}.
\end{itemize}

\section{Artigos em Periódicos}
Os exemplos abaixo ilustram referências a artigos em periódicos.
\begin{itemize}
	\item \bibentry{Hayes08}.
	\item \bibentry{Lawton08}.
\end{itemize}

\section{Artigos em Conferências}
\begin{itemize}
	\item \bibentry{Laadan10}.
	\item \bibentry{Anderson95}.
\end{itemize}

\section{Teses e Dissertações}
Seguem algumas referências a trabalhos acadêmicos, como teses, dissertações, trabalhos de conclusão de curso, etc.
\begin{itemize}
	\item \bibentry{Teixeira09}.
	\item \bibentry{Flaumann05}.
\end{itemize}

%=======================================================================
% Referências
%=======================================================================
\bibliography{exemplo}

%=======================================================================
% Exemplo de Apêndice
% O Apêndice é utilizado para apresentar material complementar elaborado
% pelo próprio autor.  Deve seguir as mesmas regras de formatação do
% corpo principal do documento.
%=======================================================================
\appendix
\chapter{Informações Complementares}

O Apêndice é o lugar para incluir textos complementares, que não são essenciais para o entendimento do assunto principal da monografia, mas que podem contribuir com informação relevante (por exemplo, uma prova matemática, uma conceituação básica, etc.).  Ele deve seguir o formato normal do documento.

%=======================================================================
% Exemplo de Anexo
% O Anexo é utilizado para a ``inclusão de materiais não elaborados pelo
% próprio autor, como cópias de artigos, manuais, folders, balancetes, etc.
% e não precisam estar em conformidade com o modelo''.
%=======================================================================
\annex
\chapter{Artigos Publicados}
Existe diferença entre os Apêndices e os Anexos.  Os apêndices trazem informação escrita pelo próprio autor do trabalho, incorporando-se ao formato da monografia como um todo.  Já um anexo é um material à parte, definido/publicado por si só, e que o autor julga conveniente ser apresentado juntamente com a monografia.  Normalmente também vai apresentar formato próprio, como um artigo publicado, um folder, uma planilha, etc.
\end{document}
